% 第九卷 6.2.1–6.2.2 节：预填充与逐步解码的时间线，以及屋顶线示意。
% (a) 块长不按时间比例；(b) 双对数坐标，X = 9.8 + 1.3 log10(I)，Y = -1.4 + log10(TFLOPS)。
% 硬件取 A100 公开规格：峰值约 312 TFLOPS（bf16），带宽约 2.0 TB/s，平衡点约 153 次每字节。
% 三个点：单请求解码 I≈1；B=32、上下文 4096 的批量解码 I≈16.5；预填充 2048 词元 I≈2000。
\begin{tikzpicture}[
  x=1cm,y=1cm,
  every node/.style={font=\figurefont\footnotesize,text=FigureInk},
  pre/.style={rectangle,rounded corners=1.8pt,draw=FigureBlue,line width=0.6pt,fill=FigureBlueFill,minimum height=7mm},
  dec/.style={rectangle,rounded corners=1.8pt,draw=FigureYellow,line width=0.6pt,fill=FigureYellowFill,minimum height=7mm,minimum width=5.2mm,inner sep=0pt},
  tok/.style={circle,draw=FigureStroke,line width=0.5pt,fill=white,minimum size=2.6mm,inner sep=0pt},
  lab/.style={font=\figurefont\scriptsize,text=FigureMuted},
  ptitle/.style={font=\figurefont\bfseries\footnotesize,text=FigureInk,anchor=west}
]
  % (a) 时间线。
  \node[ptitle] at (-0.2,1.8) {(a) 一次请求的时间线};
  \draw[-{Stealth[length=2mm]},FigureStroke,line width=0.6pt] (-0.2,-0.6) -- (8.1,-0.6) node[lab,anchor=north east] {时间};
  \node[pre,minimum width=32mm,text width=30mm,align=center,anchor=west] (P) at (0,0) {预填充\\{\scriptsize $n$个输入词元并行}};
  \foreach \k in {1,...,6} {
    \node[dec,anchor=west] (D\k) at ({3.3+(\k-1)*0.62},0) {};
  }
  \foreach \k in {1,...,6} {\node[lab] at (D\k.center) {解码};}
  \node[lab] at (7.3,0) {$\cdots$};
  % 输出词元出现的时刻。
  \node[tok] (Y1) at (P.east|-0,0.75) {};
  \foreach \k in {1,...,6} {\node[tok] at (D\k.east|-0,0.75) {};}
  \node[lab,anchor=south] at (P.east|-0,0.9) {$y_1$};
  \node[lab,anchor=south] at (D1.east|-0,0.9) {$y_2$};
  \node[lab,anchor=south] at (D2.east|-0,0.9) {$y_3$};
  % 延迟标注。
  \draw[FigureStroke,line width=0.5pt] (0,-0.45) -- (0,-0.75);
  \draw[FigureStroke,line width=0.5pt] (P.east|-0,-0.45) -- (P.east|-0,-0.75);
  \draw[{Stealth[length=1.5mm]}-{Stealth[length=1.5mm]},FigureBlue,line width=0.6pt] (0,-0.95) -- node[lab,below,text=FigureInk] {首词元延迟} (P.east|-0,-0.95);
  \draw[{Stealth[length=1.5mm]}-{Stealth[length=1.5mm]},FigureYellow,line width=0.6pt] (D3.east|-0,-0.95) -- node[lab,below,text=FigureInk] {逐词元延迟} (D4.east|-0,-0.95);
  \draw[FigureStroke,line width=0.5pt] (D3.east|-0,-0.45) -- (D3.east|-0,-0.75);
  \draw[FigureStroke,line width=0.5pt] (D4.east|-0,-0.45) -- (D4.east|-0,-0.75);

  % (b) 屋顶线。
  \node[ptitle] at (9.3,1.8) {(b) 屋顶线示意（双对数）};
  \draw[-{Stealth[length=2mm]},FigureStroke,line width=0.6pt] (9.5,-1.4) -- (14.8,-1.4);
  \draw[-{Stealth[length=2mm]},FigureStroke,line width=0.6pt] (9.5,-1.4) -- (9.5,1.45);
  \foreach \x/\t in {9.8/1,11.1/10,12.4/100,13.7/1000} {
    \draw[FigureStroke,line width=0.5pt] (\x,-1.4) -- (\x,-1.47);
    \node[lab,anchor=north] at (\x,-1.47) {\t};
  }
  \foreach \y/\t in {-0.4/10,0.6/100} {
    \draw[FigureStroke,line width=0.5pt] (9.5,\y) -- (9.43,\y);
    \node[lab,anchor=east] at (9.43,\y) {\t};
  }
  \node[lab,anchor=north] at (12.15,-1.82) {算术强度（次/字节）};
  \node[lab,anchor=south,rotate=90] at (8.95,0.0) {可达速率（TFLOPS）};
  \draw[FigureStroke,line width=1.1pt] (9.56,-1.28) -- (12.64,1.094) -- (14.6,1.094);
  \draw[FigureGrid,dashed,line width=0.5pt] (12.64,1.094) -- (12.64,-1.4);
  \node[lab,anchor=north west] at (12.72,0.5) {平衡点$\approx153$};
  \node[lab,anchor=west] at (11.7,-0.75) {带宽限制};
  \node[lab,anchor=south] at (13.5,1.12) {算力限制};
  \fill[FigureYellow] (9.8,-1.09) circle (0.075);
  \node[lab,anchor=west,text=FigureInk] at (9.88,-1.02) {单请求解码};
  \fill[FigureYellow] (11.38,0.13) circle (0.075);
  \node[lab,anchor=south east,text=FigureInk,align=right] at (11.33,0.17) {批量32解码\\（上下文4096）};
  \fill[FigureBlue] (14.09,1.094) circle (0.075);
  \node[lab,anchor=north,text=FigureInk] at (14.09,1.0) {预填充};
\end{tikzpicture}
